\documentclass[../../../main]{subfiles}

\ifSubfilesClassLoaded{%
    \zexternaldocument*{../../1.Category-Theory/exported/aux/main}%
    \zexternaldocument*{../1.Additive-and-abelian-categories/exported/aux/main}%
    \zexternaldocument*{../2.Complexes-and-cohomology/exported/aux/main}%
}{}

\begin{document}
\section{The mapping cone of a morphism}

\begin{definition}[Mapping cone] \label{def:mapping-cone}
  Let \(\alpha^{\bullet}: L^{\bullet} \to M^{\bullet}\) be a \hyperref[def:complex-morphism]{morphism of complexes}.
  We define the \textbf{mapping cone} \(\mCone{\alpha}{\bullet}\) as the \hyperref[def:complex]{complex} with objects
  \begin{align*}
    \forall i \in \mathbb{Z} \quad \mCone{\alpha}{i} = L[1]^{i} \oplus M^{i} = L^{i+1} \oplus M^{i}
  \end{align*}
  and differentials
  \begin{align*}
    \forall i \in \mathbb{Z} \quad d_{\mCone{\alpha}{\bullet}}^{i}(\ell, m) = (-d_{L^{\bullet}}^{i+1}(\ell), \alpha^{i+1}(\ell) + d_{M^{\bullet}}^{i}(m))
  \end{align*}
  Pictorially
  \begin{center}
    \begin{tikzcd}[row sep=tiny]
      \cdots \arrow[r] \arrow[ddr] & L^{i+1} \arrow[r, "-d_{L}^{i+1}"] \arrow[ddr, "\alpha^{i+1}"] & L^{i+1} \arrow[r] \arrow[ddr] & \cdots \\
      & \oplus & \oplus \\
      \cdots \arrow[r]  & M^{i} \arrow[r, "d_{M}^{i}"'] & M^{i+1} \arrow[r] & \cdots
    \end{tikzcd}
  \end{center}
\end{definition}

\begin{remark}
  We have \hyperref[def:complex-morphism]{morphisms of complexes} \(M^{\bullet} \to \mCone{\alpha}{\bullet}\) and \(\mCone{\alpha}{\bullet} \to L[1]^{\bullet}\) induced by the morphisms \(M^{i} \to L[1]^{i} \oplus M^{i} \to L[1]^{i}\).

  Furthermore, since the we have \hyperref[def:short-exact-sequence]{short exact sequences}
  \begin{center}
    \begin{tikzcd}
      0 \arrow[r] & M^{i} \arrow[r] & L[1]^{i} \oplus M^{i} \arrow[r] & L[1]^{i} \arrow[r] & 0
    \end{tikzcd}
  \end{center}
  the sequence of \hyperref[def:complex]{complexes}
  \begin{center}
    \begin{tikzcd}
      0 \arrow[r] & M^{\bullet} \arrow[r] & L[1]^{\bullet} \oplus M^{\bullet} \arrow[r] & L[1]^{\bullet} \arrow[r] & 0
    \end{tikzcd}
  \end{center}
  is \hyperref[def:short-exact-sequence]{exact}.
\end{remark}

\begin{proposition} \label{prop:connecting-iso-cohomology}
  There exists an exact triangle
  \begin{center}
    \begin{tikzcd}[column sep=tiny]
      & \cohomology{\bullet}(M^{\bullet}) \arrow[dr] & \\
      \cohomology{\bullet}(L[1]^{\bullet}) \arrow[ur, "\delta" near start, "+1" near end] & & \cohomology{\bullet}(\mCone{\alpha}{\bullet}) \arrow[ll]
    \end{tikzcd}
  \end{center}
  where the connecting morphism \(\delta\) is the \hyperref[thm:LES-in-cohomology]{morphism induced by \(\alpha^{\bullet}\) in cohomology}.
\end{proposition}
\begin{proof}
  The existence of the triangle comes directly from \autoref{thm:LES-in-cohomology}, so we only need to check that \(\delta\) agrees with the morphism induced by \(\alpha^{\bullet}\).

  By chasing the diagram
  \begin{center}
    \begin{tikzcd}
      & \vdots & \vdots & \vdots & \\
      0 \arrow[r] & M^{i} \arrow[r] \arrow[u] & L^{i+1} \oplus M^{i} \arrow[r] \arrow[u] & L^{i+1} \arrow[r] \arrow[u] & 0 \\
      0 \arrow[r] & M^{i-1} \arrow[r] \arrow[u] & L^{i} \oplus M^{i-1} \arrow[r] \arrow[u] & L^{i} \arrow[r] \arrow[u] & 0 \\
      & \vdots \arrow[u] & \vdots \arrow[u] & \vdots \arrow[u] &
    \end{tikzcd}
  \end{center}
  with a class \(\cohomology{i-1}(L[1]^{\bullet}) = \cohomology{i}(L^{\bullet})\) represented by \(\ell \in L^{i}\) (such that \(d_{L^{\bullet}}^{i}(\ell) = 0\)) we find
  \begin{center}
    \begin{tikzcd}
      \alpha^{i}(\ell) \arrow[r, mapsto] & (0, \alpha^{i}(\ell)) \\
      & (\ell, 0) \arrow[u, mapsto] \arrow[r, mapsto] & \ell
    \end{tikzcd}
  \end{center}
  which verifies the claim.
\end{proof}

\begin{corollary} \label{cor:exactness-mapping-cone}
  Let \(\alpha^{\bullet}: L^{\bullet} \to M^{\bullet}\) be a \hyperref[def:complex-morphism]{morphism of complexes}.
  Then the \hyperref[thm:LES-in-cohomology]{induced morphism} \(\cohomology{\bullet}(L^{\bullet}) \to \cohomology{\bullet}(M^{\bullet})\) is an \hyperref[def:categorical-isomorphism]{isomorphism} if and only if the \hyperref[def:mapping-cone]{mapping cone} \(\mCone{\alpha}{\bullet}\) is an \hyperref[def:exact-complex]{exact complex}.
\end{corollary}
\begin{proof}
  Using \autoref{prop:connecting-iso-cohomology}, consider the following segment of the \hyperref[def:exact-complex]{long exact sequence}
  \begin{center}
    \begin{tikzcd}
      \cdots \arrow[r] & \cohomology{i-1}(\mCone{\alpha}{\bullet}) \arrow[r] \arrow[d, phantom, ""{coordinate, name=Z}] & \cohomology{i}(L^{\bullet}) \arrow[dll, "\alpha^{\ast}", rounded corners, to path={
        -- ([xshift=2ex]\tikztostart.east)
        |- (Z) [near end]\tikztonodes
        -| ([xshift=-2ex]\tikztotarget.west)
        -- (\tikztotarget)}] \\
      \cohomology{i}(M^{\bullet}) \arrow[r] & \cohomology{i}(\mCone{\alpha}{\bullet}) \arrow[d, phantom, ""{coordinate, name=W}] \arrow[r] & \cohomology{i+1}(L^{\bullet}) \arrow[dll, "\alpha^{\ast}", rounded corners, to path={
        -- ([xshift=2ex]\tikztostart.east)
        |- (W) [near end]\tikztonodes
        -| ([xshift=-2ex]\tikztotarget.west)
        -- (\tikztotarget)}] \\
      \cohomology{i+1}(M^{\bullet}) \arrow[r] & \cohomology{i+1}(\mCone{\alpha}{\bullet}) \arrow[r] & \cdots
    \end{tikzcd}
  \end{center}
  \begin{itemize}
    \item[\(\implies\)] Suppose \(\alpha^{\ast}\) is an \hyperref[def:categorical-isomorphism]{isomorphism}.
          Since \hyperref[thm:complex-category-is-abelian]{\(\Complex{\category{A}}\) is abelian} by \autoref{lemma:mono-and-epi-is-iso} \(\alpha^{\ast}\) is both a \hyperref[def:mono-epi-morphism]{monomorphism} and an \hyperref[def:mono-epi-morphism]{epimorphism}. Consider
          \begin{center}
            \begin{tikzcd}
              \cdots \arrow[r, "\alpha^{\ast}"] & \cohomology{i}(M^{\bullet}) \arrow[r, "f"] & \cohomology{i}(\mCone{\alpha}{\bullet}) \arrow[r, "g"] & \cohomology{i+1}(L^{\bullet}) \arrow[r, "\alpha^{\ast}"] & \cdots
            \end{tikzcd}
          \end{center}
          Using that the first \(\alpha^{\ast}\) is a surjective and \hyperref[def:exact-complex-sequence]{exactness} we have that
          \begin{align*}
            \ker f = \im \alpha^{\ast} = \cohomology{i}(M^{\bullet})
          \end{align*}
          i.e: \(f\) is the zero map.

          Using that the last \(\alpha^{\ast}\) is injective and \hyperref[def:exact-complex-sequence]{exactness} we have that
          \begin{align*}
            \ker g = \ker \alpha^{\ast} = \{0\}
          \end{align*}
          i.e: \(g\) is the zero map.

          Therefore we have that \(\cohomology{i}{\mCone{\alpha}{\bullet}} = 0\), meaning that the \hyperref[def:complex]{complex} is \hyperref[def:exact-complex]{exact}.
    \item[\(\impliedby\)] Suppose \(\mCone{\alpha}{\bullet}\) is an \hyperref[def:exact-complex]{exact complex}.
          Therefore since its \hyperref[def:cohomology]{cohomology} vanishes at every index the sequence becomes
          \begin{center}
            \begin{tikzcd}
              \cdots \arrow[r] & 0 \arrow[r] & \cohomology{i}(L^{\bullet}) \arrow[r, "\alpha^{\ast}"] & \cohomology{i}(M^{\bullet}) \arrow[r] & 0 \arrow[r] & \cdots
            \end{tikzcd}
          \end{center}
          Therefore by \hyperref[def:exact-complex-sequence]{exactness} we have that \(\alpha^{\ast}\) is an \hyperref[def:categorical-isomorphism]{isomorphism}.
  \end{itemize}
\end{proof}

\end{document}
